% BHU PYQ -- Auto-generated & error-corrected
\documentclass[11pt,a4paper]{article}

\usepackage[a4paper,top=2.5cm,bottom=2.5cm,left=2.8cm,right=2.8cm,headheight=14pt]{geometry}
\usepackage{amsmath,amssymb}
\usepackage[T1]{fontenc}
\usepackage[utf8]{inputenc}
\usepackage{lmodern}
\usepackage{microtype}
\usepackage{fancyhdr}
\usepackage{enumitem}
\usepackage{hyperref}
\usepackage{lastpage}
\usepackage{parskip}

\hypersetup{colorlinks=true,urlcolor=black,linkcolor=black,
  pdftitle={STA/STB-603: Elements of Stochastic Processes -- BHU PYQ},pdfauthor={Banaras Hindu University}}

\pagestyle{fancy}\fancyhf{}
\renewcommand{\headrulewidth}{0.4pt}\renewcommand{\footrulewidth}{0.4pt}
\fancyhead[L]{\small   \textit{Banaras Hindu University}}
\fancyhead[R]{\small   \textit{STA/STB-603: Elements of Stochastic Processes}}
\fancyfoot[C]{\small Page  \thepage\ of \pageref{LastPage}}


\newlist{parts}{enumerate}{2}
\setlist[parts,1]{label=(\alph*),leftmargin=*,itemsep=5pt,topsep=3pt}
\setlist[parts,2]{label=(\roman*),leftmargin=*,itemsep=3pt,topsep=2pt}

\newcommand{\pts}[1]{\hfill   \textit{[#1]}}


\begin{document}


\begin{center}
  {\large\bfseries BANARAS HINDU UNIVERSITY}\\[2pt]
  {
\normalsize B.A./B.Sc. (Hons.) Semester VI Examination, 2013-14}\\[6pt]
  \rule{0.8\linewidth}{0.5pt}\\[6pt]
  {\large\bfseries Paper No. STA/STB-603: Elements of Stochastic Processes}\\[4pt]
  {
\normalsize Time: 3 Hours\hfill Maximum Marks: 70}
\end{center}
\rule{\linewidth}{0.4pt}
\smallskip



\noindent   \textit{Instructions: Answer any   \textbf{five} questions. All questions carry equal marks. Symbols have their usual meanings.}
\bigskip

\medskip


\noindent   \textbf{Question 1.}

What do you mean by a stochastic process? Discuss its classification according to state space and time domain. Give one example of each type. \pts{14}

\medskip\rule{\linewidth}{0.2pt}

\medskip


\noindent   \textbf{Question 2.}

Describe the Gambler's Ruin Problem. Obtain the expression for the expected duration of the game for $p = \frac{1}{2}$ and for $p 
eq \frac{1}{2}$. \pts{14}

\medskip\rule{\linewidth}{0.2pt}

\medskip


\noindent   \textbf{Question 3.}

\begin{parts}
  \item Define a Markov chain. Give two examples. \pts{7}
  \item Define the $n$-step transition probability. Also prove the Chapman--Kolmogorov equation:
  \[
    P_{ij}^{(n+m)} = \sum_{k} P_{ik}^{(n)} P_{kj}^{(m)}.
  \]\pts{7}
\end{parts}

\medskip\rule{\linewidth}{0.2pt}

\medskip


\noindent   \textbf{Question 4.}

\begin{parts}
  \item Define persistent, transient, ergodic, and essential states of a Markov chain. \pts{7}
  \item Describe stationary distribution for a Markov chain. Give conditions under which it exists. \pts{7}
\end{parts}

\medskip\rule{\linewidth}{0.2pt}

\medskip


\noindent   \textbf{Question 5.}

\begin{parts}
  \item Describe the time-dependent Poisson process. Obtain the expression for $P_n(t)$, mean and variance. \pts{7}
  \item If $\{X(t),\, t \geq 0\}$ is a Poisson process, obtain $P[X(s) = k \mid X(t) = n]$ for $s < t$. \pts{7}
\end{parts}

\medskip\rule{\linewidth}{0.2pt}

\medskip


\noindent   \textbf{Question 6.}

Describe the pure death process. Obtain the expression for $P_n(t)$ in the case of a linear death process. \pts{14}

\medskip\rule{\linewidth}{0.2pt}

\medskip


\noindent   \textbf{Question 7.}

Describe the linear birth--death process. Derive the expression for its probability generating function. \pts{14}

\medskip\rule{\linewidth}{0.2pt}

\medskip


\noindent   \textbf{Question 8.}

Describe the Galton--Watson branching process. For this process, prove that the probability generating function satisfies $P_n(s) = P_{n-1}(P(s)) = P(P_{n-1}(s))$, where $P_n(s) = \sum_{k=0}^{\infty} P(Z_n = k)\,s^k$. \pts{14}

\vfill
\rule{\linewidth}{0.4pt}

\begin{center}\small
  \href{https://bhu-exam.vercel.app/}{Click here to visit our website}\\[4pt]
  \href{https://me-aryan.vercel.app/}{Click here to visit our contributor}\\[4pt]
  \href{https://quashan-me.vercel.app/}{Click here to visit our contributor}
\end{center}

\end{document}